\diarytitle{0119}{極座標変換を用いた四分円板上の広義重積分の計算}

極座標変換 $x = r\cos\theta,\ y = r\sin\theta$ を用いる。ヤコビ行列式は $\left|\dfrac{\partial(x,y)}{\partial(r,\theta)}\right| = r$ である（問題0117の結果と同じ計算による）。領域 $D$ は第1象限内の半径 $2$ の円板の内部であるから、$(x,y)$ 平面の $D$ は $(r,\theta)$ 平面の
\[
D' = \{(r,\theta) \mid 0 \le r \le 2,\ 0 \le \theta \le \tfrac{\pi}{2}\}
\]
に対応する。また $x^{2}+y^{2} = r^{2}$ より被積分関数は $\dfrac{1}{\sqrt{4-r^{2}}}$ となるから、
\[
I = \int_{0}^{\pi/2} \int_{0}^{2} \frac{1}{\sqrt{4-r^{2}}} \cdot r\, dr\, d\theta
\]
$r$ に関する積分は、$t = 4 - r^{2}$ とおくと $dt = -2r\, dr$ より
\[
\int_{0}^{2} \frac{r}{\sqrt{4-r^{2}}}\, dr
= \left[ -\sqrt{4-r^{2}} \right]_{0}^{2}
= -0 - (-2) = 2
\]
（$r \to 2^{-}$ でも被積分関数は可積分特異点であり、この広義積分は収束する。）よって
\[
I = \int_{0}^{\pi/2} 2\, d\theta = 2 \cdot \frac{\pi}{2} = \boxed{\; \pi \;}
\]
（検算: 極座標での数値積分によっても $I = \pi \approx 3.1416$ と一致することを確認できる。）
